\chapter{Conclusion and Perspectives}
\label{chap:conc}
\graphicspath{Conclusion/img/}


\section{Conclusion}

This study investigated how a multi-layered NLP framework could identify,
analyse, and explain differences between AI-in-education policy and stakeholder
evidence surrounding its implementation. Rather than treating the policy--practice
gap as a measurable quantity, the study examined one observable layer of that
gap through policy and sentiment discourse. The framework combined corpus
construction, topic modelling, semantic causal-claim coverage, causal-frame networks,
synthetic-data sensitivity analysis, and exploratory LLM-assisted interpretation.

The topic-modelling stage established the thematic structure required for subsequent
analysis. LDA, BERTopic-style clustering, and NMF provided complementary views of the
corpora, while NMF supplied the reference topic space used by the causal analyses.
These methods showed that policy and sentiment frequently addressed comparable
AI-in-education concerns, but thematic presence alone did not establish equivalent
interpretation or emphasis. This motivated the analysis of causal language and the
relationships connecting educational concerns.

The two causal NLP methods provided the principal evidence for a selective
policy--sentiment causal gap. Semantic coverage showed substantial correspondence
between the corpora alongside persistent differences in the availability of close
causal-claim counterparts. Causal-frame analysis further showed that comparable
subjects could occupy different causal roles. Policy placed greater emphasis on
enabling mechanisms, institutional requirements, and intended improvements, whereas
sentiment more often foregrounded expressed causes and consequences. The strongest
differences concerned learning and assessment, curriculum development and oversight,
and the organisation of institutional mechanisms around these concerns. The findings
therefore indicated structured differences rather than a uniform separation between
policy and sentiment.

The LLM-assisted stage provided a complementary interpretation of this evidence.
It identified causal or conditional relationships that could remain outside fixed
cue extraction and provided explicit explanations for evidence relationships that
semantic similarity represented numerically. It also exposed cases where superficial
cue matching produced questionable causal candidates. These observations did not
establish LLM superiority. Instead, they illustrated the benefit of combining
reproducible corpus-level NLP measurements with evidence-grounded interpretation
while retaining human review.

Overall, the research question was addressed by demonstrating that a multi-layered
NLP framework can locate thematic correspondence, measure differences in causal-claim
coverage, reveal structural differences in causal framing, and connect these patterns
back to traceable sentence-level evidence. Its principal contribution lies in
distinguishing different forms of policy--sentiment divergence rather than reducing
them to a single gap score. The resulting evidence concerns expressed causal
discourse and provides a computational view of one dimension of the broader
policy--practice problem; it does not demonstrate real-world causal effects,
implementation outcomes, or policy effectiveness.

\section{Limitations}

Three limitations are most relevant to the interpretation of the findings. First,
policy and sentiment evidence were unevenly distributed across countries, with
substantially less original sentiment evidence available at national level.
This constrained country-level comparison, particularly for Ireland and France.
The imbalance also limited direct comparison of policy priorities with stakeholder
responses within the same national contexts. Moreover, stakeholder discourse captures
perceptions, concerns, and experiences rather than direct classroom implementation.
The available evidence therefore reflects how policy is received and discussed,
but not how consistently it is enacted in educational practice. The
policy--sentiment gap consequently represents one observable dimension of the broader
policy--practice gap.

Second, the causal NLP results remain dependent on modelling choices. Semantic
coverage relies on predefined causal cues, sentence embeddings, and
nearest-neighbour similarity, while causal-frame analysis depends on relation-family
definitions, topic projection, and fallback rules. Less formulaic causal relationships may therefore
remain outside the extracted claim inventory, and some frame assignments required
document-topic fallback. Robustness analyses showed that the main structural
differences persisted across alternative specifications, but these methodological
dependencies remain.

Third, the LLM-assisted analysis was exploratory and operated on sampled evidence
rather than the complete corpus. Its findings relied on automated analyst and
verifier judgements, without a completed human-labelled evaluation of precision,
recall, or classification accuracy. The LLM results should therefore be treated as
evidence-grounded interpretations requiring further human validation. Synthetic
sentiment was similarly restricted to sensitivity testing and did not contribute to
the empirical conclusions.

\section{Perspectives}

Future research could first expand the empirical evidence base. Larger and more
balanced stakeholder corpora could include teachers, students, parents, school
leaders, and policy implementers across additional countries. Longitudinal evidence
would also make it possible to examine whether policy--sentiment relationships change
as AI guidance, regulation, technologies, and educational practices evolve. Direct
implementation evidence, including institutional reports, interviews, surveys, and
classroom studies, would provide an important bridge from the policy--sentiment layer
examined here toward the broader policy--practice gap.

A second direction concerns methodological extension. Causal extraction could move
beyond a fixed cue inventory through multilingual contextual models with improved
cross-lingual semantic generalisation across diverse contexts while preserving the
provenance and auditability provided by the current pipeline. Human annotation of
causal claims and frame relationships would enable systematic evaluation of
extraction precision and recall. Topic and causal-frame assignments could likewise
be assessed against manually validated reference sets, strengthening the empirical
basis for comparing modelling choices.

The LLM-assisted component also offers a clear extension path. Applying the
analyst--verifier procedure across the complete corpus, followed by structured human
evaluation, would allow its contribution to be assessed more rigorously against the
causal NLP methods. Future work could examine where LLM interpretation consistently
adds evidence unavailable to fixed extraction, where it introduces unsupported
inferences, and how human review can be incorporated efficiently without weakening
traceability or reproducibility.

More broadly, the framework could support continuous comparison between evolving AI
policy and stakeholder evidence. Rather than treating policy analysis as a static
snapshot, future systems could identify emerging concerns, changes in causal framing,
and areas where institutional guidance no longer corresponds closely to reported
experience. Combined with representative empirical data and human validation, this
would extend the present framework from retrospective discourse analysis toward a
more systematic instrument for studying how AI-in-education governance develops in
practice.